\documentclass[../main.tex]{subfiles}

\begin{document}

\ifSubfilesClassLoaded{

    \tableofcontents
    
}{}

\chapter{ Summer Exercises }

\section{Continuous probability space}

\begin{exercise}{Bertrand's Paradox}{}
    Show that the probability of the event: a randomly chosen chord of the unit circle passing at a distance of $\frac{1}{2}$ from the center, is $\frac{1}{3}$ or $\frac{1}{2}$ depending on the chosen model. Compare the two models.
\end{exercise}

\begin{proof}
    \begin{enumerate}
        \item Let \(   \Omega   \) be the collection of the pair \(  \left( l, d \right)   \), where \(  l\in \left[ 0,2 \pi \right]   \), \(  d \in \left[ 0,1 \right]   \).  Let \(  \mathcal{F}= \mathcal{P}\left(  \Omega  \right)   \).  By randomly chosing a chord of the unit circle, we mean randomly chosing a \(  \left( l,d \right) \in \left[ 0,2 \pi \right]\times \left[ 0,1 \right]     \) , then we identify the \(  \left( l,d \right)   \) with the chord whose radius is  \(  l  \) and the distance between the center point  of the chord with that of circle is \(  d  \). Let \(  P  \) be the normalized  Lebesgue measure on \(  \left[ 0,2 \pi \right]\times \left[ 0,1 \right]    \), then the probability we discuss is \(  \frac{1}{2}  \).     
        \item Let \(   \Omega ^{\prime}   \) be all the    points on the boundary of the unit circle.Let \(  \mathcal{F}^{\prime} = \mathcal{P}\left(  \Omega ^{\prime}  \right)   \).   First , fix a point \(  p  \) on \(   \Omega ^{\prime}   \).     By chosing a chord of the unit circle, we mean randomly chosing a point  \(q\in    \Omega ^{\prime}   \), and identify the point with the chord connecting \(  p  \)   and \(  q  \). Since  the chord is \(  \frac{1}{2}  \) passing at the center if and only if  the angle (which is range from \(  0  \) to \(    \pi  \)  ) corrospoing to the  chord is greater than \(  \frac{ \pi }{3 }   \). Similarly let \(  P  \) be the normalized  Lebesgue measure on \(  \left(  \Omega ^{\prime} ,\mathcal{F}^{\prime}  \right)   \). Then the probability we discuss is \(  \frac{1}{3}  \) .     
    \end{enumerate}
    
\end{proof}

\begin{exercise}{Change of Variables}{}
    
\begin{enumerate}
    \item Let $X$ be a uniform random variable on $]-\frac{\pi}{2}, \frac{\pi}{2}[$. What is the distribution of $\tan X$?
    \item Let $X$ be a random variable with the density $f(x) = \frac{\mathbf{1}_{[0,1]}(x)}{\log 2(1 + x)}$ show that $\frac{1}{X} - \lfloor \frac{1}{X} \rfloor$ has the same distribution as $X$.
\end{enumerate}
\end{exercise}
\begin{proof}
    \begin{enumerate}
        \item Denote \(  \tan  X  \) by \(  Y  \).  Then \[
        F_{Y}\left( x \right)= P\left( \tan X\le x \right)  = P\left( X\le \arctan x \right) 
        \]   \[
        F_{Y}\left( x \right)=  \int_{-\infty}^{\arctan x} \frac{1 }{ \pi } 1_{\left( \frac{- \pi }{2 }, \frac{ \pi }{2 }   \right) }\,\mathrm{d} x= \frac{1 }{ \pi } \arctan x+ \frac{1 }{2 }  
        \]the density function of \(  F_{Y}  \) is \[
        f_{Y}\left( x \right)= \frac{1 }{ \pi }\frac{1 }{x^{2}+ 1 }   
        \] 
        \item When \(  x> 0  \),   \[
        P\left( \frac{1 }{X }  \le x\right)=  P \left( X\ge \frac{1 }{x }  \right) =  \int_{-\infty}^{\frac{1 }{x } } 1_{\left[ 0,1 \right] } \frac{1 }{\log 2\left( 1+y \right)  }\,\mathrm{d} y=  
        \]
    \end{enumerate}
    
\end{proof}

\begin{exercise}{}{}
    Determine the expectation and variance of a normal distribution, gamma distribution, and log-normal
 distribution
\end{exercise}
\begin{proof}
    \begin{enumerate}
        \item If \(  X \sim  \mathcal{N}\left( \mu , \sigma ^{2} \right)  \) , then the PDF is  \[
        f\left( x \right)= \frac{1 }{\sqrt{2 \pi \sigma ^{2}} }\exp \left( - \frac{\left( x-\mu  \right)^{2}  }{ 2 \sigma ^{2}}  \right)   
        \] Let \(  y= \frac{x-\mu  }{ \sigma  }   \) , then \(  x =  \sigma y+ \mu   \), \[
        \begin{aligned}
        E\left[ X \right]&= \frac{1 }{\sqrt{2 \pi \sigma ^{2}} }  \sigma  \int_{\mathbb{R} }\left(  \sigma y+ \mu  \right)\exp \left( -\frac{1}{2}y^{2} \right) \,\mathrm{d} y     \\ 
         &=  \frac{ \sigma  }{ \sqrt{2 \pi} }\int_{\mathbb{R} }\exp \left( -\frac{1}{2}y^{2} \right)\,\mathrm{d} \left( \frac{1}{2}y^{2} \right)+  \frac{\mu  }{ \sqrt{2 \pi} } \int_{\mathbb{R} }   \exp \left( -\frac{1 }{2 }y^{2}  \right)\,\mathrm{d} y  \\ 
          &= 0+  \mu = \mu  
        \end{aligned}
        \] To caculate \(  \mathrm{Var}\left( X \right)   \), we assume that   \(  \mu = 0  \). Then \[
        \begin{aligned}
        \mathrm{Var}\left( X \right)&= E\left[ \left( X-\mu  \right)^{2}  \right] = E\left[ X^{2} \right]   \\ 
         &= E\left[ X^{2} \right] =  \int_{S} \left( X\left( x \right)  \right)^{2} \,\mathrm{d} P\\ 
          &=  \int_{\mathbb{R} }x^{2} \,\mathrm{d} P_{X}\\ 
           &=  \int_{\mathbb{R} } x^{2}f\left( x \right)\,\mathrm{d} x =  \int_{-\infty}^{\infty}x^{2}\frac{1 }{\sqrt{2 \pi \sigma ^{2}} }\exp \left( -\frac{x^{2} }{2 \sigma ^{2} }  \right)\,\mathrm{d} x    
        \end{aligned}
        \] Let \(  z = \frac{x }{ \sigma  }   \), then \[
        \begin{aligned}
        \mathrm{Var}\left( X \right)&=  \int_{-\infty}^{\infty} \sigma ^{2}z^{2} \frac{1 }{\sqrt{2 \pi} \sigma  }\exp \left( -\frac{1}{2}z^{2} \right) \sigma \,\mathrm{d} z\\ 
         &= \sigma ^{2} \int_{-\infty}^{\infty}z^{2}\frac{1 }{\sqrt{2 \pi} } \exp \left( -\frac{1 }{2 }z^{2}  \right)\,\mathrm{d} z\\ 
          &=   \sigma ^{2} \mathrm{Var}\left( Z\sim \mathcal{N}\left( 0,1 \right)  \right)\\ 
           &=  \sigma ^{2}        
        \end{aligned}
        \] Specifically,  consider \[
        I\left( a \right)= \int_{-\infty}^{\infty}\exp \left( -a\frac{z^{2} }{2 }  \right)\,\mathrm{d} z = \sqrt{\frac{2 \pi }{a } }  
        \] and take differention of \(  a  \).  
        \item A Gamma distribution has the density fuction \[
        f\left( x;\alpha ,\beta  \right)= \frac{\beta ^{\alpha } }{ \Gamma \left(  \alpha  \right)  }x^{\alpha -1}  e^{-\beta x},\quad x> 0, \alpha , \beta > 0
        \]
        \begin{note}
            事件以平均恒定的速率 \(  \beta   \)发生,等待第一次时间发生的事件服从指数分布, 等待第 \(  \alpha   \)次时间发生的时间服从 \(   \Gamma \left( \alpha ,\beta  \right)   \)分布.   
        \end{note}
        \[
        E[X]= \frac{\alpha  }{\beta  } ,\quad \mathrm{Var}\left( X \right)= \frac{ \alpha  }{\beta ^{2} }  
        \]
        \item by a log-normal distribution \(  X  \), we mean \(  Y= \log X  \) is a normal distribution. Then \[
      F_{X}\left( x \right)=   P\left( X\le x \right)= P\left( Y\le \log x \right)= F_{Y}\left( \log x \right)  ,\quad x> 0
        \] \[
        f_{X}=  \frac{\mathrm{d}F_{X}}{\mathrm{d}x}= f_{Y}\left( \log x \right)\frac{1 }{x }  = \frac{1 }{\sqrt{2 \pi  \sigma ^{2}} }\exp \left( - \frac{\left( \log x- \mu  \right)^{2}  }{2 \sigma ^{2} }  \right)  \frac{1 }{x } ,\quad x> 0
        \]
    \end{enumerate}
    
\end{proof}

\begin{exercise}{Couples of Random Variables}{}
    \begin{enumerate}
\item Let $X$ and $Y$ be independent random variables with distribution $\mathcal{N}(0, 1)$. Determine the distribution of the pair $(X/Y, Y)$, then that of $Y$. Are $X/Y$ and $Y$ independent?
\end{enumerate}
\end{exercise}
\begin{proof}
    \[
    f_{X}= \frac{1 }{\sqrt{2 \pi} } \exp \left( - \frac{x^{2} }{2 }  \right),\quad f_{Y}= \frac{1 }{\sqrt{2 \pi} }\exp \left( -\frac{y^{2} }{2 }  \right)    
    \] Let \(  U = X/Y, V= Y  \), then \(  X = UV, Y = V  \), \[
    f_{U,V}\left( u,v \right)= f_{X,Y}\left( x\left( u,v \right),y\left( u,v \right)   \right) \left|  \det \begin{pmatrix} 
        \frac{\partial x}{\partial u}& \frac{\partial x}{\partial v}\\ 
         \frac{\partial y}{\partial u}&\frac{\partial y}{\partial v} 
    \end{pmatrix}  \right| 
    \]  where  \[
    f_{X,Y}\left( x,y \right)= f_{X}\left( x \right)f_{Y}\left( y \right)= \frac{1 }{2 \pi } \exp \left( -\frac{x^{2}+ y^{2} }{2 }  \right)     
    \] Then \[
    f_{U,V}\left( u,v \right)= \frac{1 }{2 \pi }\exp \left( -\frac{u^{2}v^{2}+ v^{2} }{2 }  \right)\left| \det \begin{pmatrix} 
        v&u\\ 
         0&1 
    \end{pmatrix}  \right|    =  \frac{1 }{2 \pi }\exp \left( -\frac{v^{2}\left( u^{2}+ 1 \right)  }{2 }  \right)\left| v \right|   
    \] \[
    f_{Y}\left( y \right)= f_{V}\left( v \right)= \frac{1 }{\sqrt{2 \pi} }\exp \left( -\frac{v^{2} }{2 }  \right)    
    \] \[
 \left(    f_{U,V}/ f_{V} \right) \left( u,v \right) = \frac{\left| v \right|  }{\sqrt{2 \pi} } \exp \left(-\frac{u^{2}v^{2} }{2 }   \right)  
    \] is not independent with \(  u  \). Thus \[
    f_{U,V}\neq f_{U}f_{V}
    \] \(  \frac{X}{Y}  \) and \(  Y  \) are not independent.  
\end{proof}
\begin{exercise}{ Estimation of a Support Paramete}{}
    
Let $(X_n)$ be a sequence of i.i.d. random variables uniformly distributed on $[0,\theta]$ ($\theta > 0$) and $M_k = \sup_{1 \le n \le k} X_n$.
\begin{enumerate}
\item Calculate the density of $M_k$.
\item Calculate the expectation and variance of $M_k$.
\item We propose to estimate the parameter $\theta$. What do you think of $M_k$ as an estimator, is it unbiased? Propose an unbiased estimator $\hat{\theta}_k$ constructed from $M_k$. Show that this estimator converges in probability towards $\theta$ and compare it to the unbiased estimator constructed from the empirical mean of the observations.
\item Show that $k\left(1 - \frac{M_k}{\theta}\right)$ converges in distribution.
\end{enumerate}
\end{exercise}

\begin{proof}
  \begin{enumerate}
    \item   \[
    M_{k}\le m\iff  \forall n, X_{n}\le m
    \] \[
    P\left( M_{k}\le m \right)= \prod _{n = 1}^{k}P\left( X_{n}\le m \right)  
    \] \[
    F_{M_{k}}=  \left( F_{X} \right)^{n} = \left( 1_{\left[ 0, \theta  \right] }\frac{x }{ \theta  }  \right)^{n}= 1_{\left[ 0, \theta  \right] }\frac{x^{n} }{ \theta ^{n} }   
    \] \[
    f_{M_{k}}= \frac{\mathrm{d}F_{M_{k}}}{\mathrm{d}x}=  1_{\left[ 0, \theta  \right] } k\frac{x^{k-1} }{  \theta ^{k}} 
    \]
    \item  \[
    E\left[ M_{k} \right]= \int_{\mathbb{R} } x f_{M_{k}} \left( x \right)\,\mathrm{d} x=   \int_{0}^{ \theta }k\frac{x^{k} }{ \theta ^{k} }\,\mathrm{d} x = \frac{k }{ \theta ^{k} }\frac{1 }{k+ 1 } \theta ^{k+ 1}= \frac{k \theta  }{k+ 1 }   
    \] \[
 \begin{aligned}
    \mathrm{Var}\left( M_{k} \right)= \int_{\mathbb{R} }x^{2}f_{M_{k}}\left( x \right)\,\mathrm{d}-\mu ^{2} x&=  \int_{0}^{ \theta }k\frac{x^{k+ 1} }{ \theta ^{k} }\,\mathrm{d} x- \mu ^{2}\\ 
     &= \frac{k  \theta ^{2} }{k+ 2 } \\ 
      &  =  \theta ^{2}\left( \frac{k }{k+ 2 }-\frac{k^{2} }{\left( k+ 1 \right)^{2}  }   \right)  \\ 
     &=  \theta ^{2}\left( \frac{k\left( k+ 1 \right)^{2}-k^{2}\left( k+ 2 \right)   }{ \left( k+ 1 \right)^{2}\left( k+ 2 \right)  }  \right)\\ 
      &=  \theta ^{2}\frac{k }{\left( k+ 1 \right)^{2}\left( k+ 2 \right)   }  
 \end{aligned}
    \]
    \item \[
    E\left[ M_{k} \right]- \theta  \neq 0 
    \] it is not unbiased. Let \[
    \hat{M}_{k}= \frac{k+ 1 }{k }M_{k} 
    \] Then \[
    \hat{ \theta }_{k}=  \theta 
    \] By chebyshev's inequality \[
    a^{2}E\left( \left| X \right|> a  \right) \le EX^{2}
    \]only need tow show that  \[
    \lim_{k\to \infty} \mathrm{Var}\left( \hat{M}_{k} \right)
    \]
  \end{enumerate}
  
\end{proof}
\end{document}